% ===========================================================================
\appendix
\part{Reference}
% ===========================================================================

\chapter{Command reference}

\section{Daily loop}

\begin{lstlisting}[style=sh]
docker compose up -d db                  # the database must be running first
cargo leptos watch                       # builds both programs, serves on :3000, reloads
cargo leptos watch -vv                   # ... and shows the sass/tailwind command lines
\end{lstlisting}

\section{The two-build check --- run both, always}

\begin{lstlisting}[style=sh]
cargo check --no-default-features --features ssr
cargo check --lib --no-default-features --features hydrate --target wasm32-unknown-unknown

# NOT this: it succeeds while checking neither program (see Chapter 7)
# cargo check
\end{lstlisting}

\section{Migrations}

\begin{lstlisting}[style=sh]
diesel migration generate create_talks   # writes an up.sql / down.sql pair
diesel migration run                     # apply pending
diesel migration revert                  # roll back the most recent
diesel migration redo                    # revert + re-apply: THE step that proves down.sql works
diesel migration redo -a                 # ... for every migration
diesel migration list                    # what is applied
diesel database reset                    # drop, create, run all -- destroys data
diesel print-schema > src/schema.rs      # regenerate; commit with the migration
diesel print-schema | diff -u src/schema.rs -    # is schema.rs stale?
\end{lstlisting}

\section{Database}

\begin{lstlisting}[style=sh]
psql "$DATABASE_URL"                     # a shell
psql "$DATABASE_URL" -c '\dt'            # list tables
psql "$DATABASE_URL" -c '\d+ tags'       # one table: types, defaults, indexes, constraints
psql "$DATABASE_URL" -c '\ds'            # list sequences -- how defect D1 was found
psql "$DATABASE_URL" -f seeds/dev_seed.sql
\end{lstlisting}

\section{Release and run}

\begin{lstlisting}[style=sh]
cargo leptos build --release

# the binary needs its environment -- see Chapter 7
LEPTOS_OUTPUT_NAME=egonik-site \
LEPTOS_SITE_ROOT=target/site \
LEPTOS_ENV=PROD \
./target/release/egonik-site
\end{lstlisting}

\section{Diagnostics}

\begin{lstlisting}[style=sh]
# what is actually in the browser bundle?
cargo tree --target wasm32-unknown-unknown --no-default-features --features hydrate --depth 1

# why is this crate in the wasm build?
cargo tree --target wasm32-unknown-unknown --no-default-features --features hydrate -i dotenvy

# did hydration happen?
curl -s -o /dev/null -w '%{http_code}\n' localhost:3000/pkg/egonik-site.js   # must be 200

# is the page really server-rendered, or an empty shell?
curl -s localhost:3000/publications | grep -c '<li'

# does a missing page really 404?
curl -s -o /dev/null -w '%{http_code}\n' localhost:3000/nonexistent

# is a file in the build at all?  (put a syntax error in it; silence means dead)
cargo check --no-default-features --features ssr
\end{lstlisting}

\section{This document}

\begin{lstlisting}[style=sh]
cd docs && tectonic -X compile main.tex        # produces main.pdf
\end{lstlisting}

% ===========================================================================
\chapter{Glossary}

\begin{description}
  \item[SSR --- server-side rendering] Running your components on the server to produce HTML. In
    this project it is a Cargo feature, \code{ssr}, not a runtime mode.

  \item[Hydration] The browser taking over server-rendered HTML: attaching event listeners and
    reconnecting reactivity without re-creating the DOM. Enabled by the \code{hydrate} feature. If
    it fails, the page looks correct and does nothing --- Chapter~\ref{ch:runcontract}.

  \item[CSR --- client-side rendering] Shipping an empty shell and building the page entirely in the
    browser. The \code{csr} feature. Not used here, and its binary target does not currently
    compile.

  \item[Server function] An \code{async fn} marked \code{\#[server]}. The macro compiles the body
    into the server only, and replaces it on the client with an HTTP call. It is a
    \textbf{public endpoint}, not private RPC.

  \item[DTO --- data transfer object] A plain serde-serialisable struct that crosses the network.
    Distinct from the Diesel row struct by design --- Chapter~\ref{ch:dto}.

  \item[Resource] Leptos's async-data primitive for \emph{reads}. Runs during SSR and serialises its
    result into the HTML so hydration does not refetch. Contrast \code{LocalResource}, which does
    not run during SSR at all.

  \item[Action] Leptos's async primitive for \emph{writes}. Runs when dispatched, never during
    server rendering.

  \item[Suspense / Transition] Boundaries that render a fallback while nested resources load.
    \code{Suspense} shows the fallback on every reload; \code{Transition} only on the first.

  \item[Hydration mismatch] The server's markup and the client's first render disagreeing, usually
    from something non-deterministic in a \code{view!}. Produces inexplicable behaviour rather than
    an error.

  \item[r2d2 pool] A connection pool. \code{Arc}-backed, so \code{Clone} is cheap and
    \code{\&self} suffices; never wrap it in another \code{Arc}.

  \item[Blocking the runtime] Calling synchronous code (all of Diesel) inside an async task, holding
    a worker thread that cannot then serve anyone else. Fixed with \code{web::block} ---
    Chapter~\ref{ch:blocking}.

  \item[Composition root] The one place that knows how everything is wired: \file{main.rs}. Config
    is loaded there, the pool is built there, nothing below it reads the environment.

  \item[Migration] A reversible \emph{pair} of SQL scripts. If \code{down.sql} does not work, it is
    not a migration --- Chapter~\ref{ch:migrations}.
\end{description}

% ===========================================================================
\chapter{Sources}

Claims in this document come from one of three places, and it is worth being able to tell which.

\begin{description}
  \item[Observation.] Anything in a \emph{Verified} box was run on this machine against this
    repository, this database, and the installed toolchain (rustc 1.93.1, cargo-leptos 0.3.6,
    PostgreSQL 17 on port 5439). Quoted output is real, occasionally trimmed for width.

  \item[Reading the source.] Behaviour of cargo-leptos 0.3.6 (the Tailwind pipeline, the hardcoded
    \code{v4.2.1}, the CSS concatenation order, the ``Done'' substring check) and of
    \code{leptos\_actix} 0.8.7 (that \code{leptos\_routes} registers server-function paths itself,
    that \code{extract} is \code{use\_context::<Request>()} plus \code{T::extract}) was taken from
    the published crate sources rather than from prose documentation.

  \item[Documentation.] The rest: the Leptos book (\code{book.leptos.dev}), \code{docs.rs} for
    \code{leptos} 0.8.20, \code{leptos\_actix} 0.8.7, \code{server\_fn} 0.8.13, \code{diesel} 2.3,
    \code{actix-web} 4, \code{r2d2}, \code{tokio}; the Actix databases guide; the Tailwind v4
    source-detection documentation; the Rust API Guidelines; and the official
    \code{leptos-rs/start-actix} and \code{start-axum-workspace} templates.
\end{description}

Version numbers current as of \textbf{4 August 2026}. The three that will matter when they change:
\code{leptos} 0.8.20 (0.9 is in beta and will move \code{Resource} and error APIs),
\code{cargo-leptos} 0.3.6 (its pinned Tailwind version is a build-behaviour change), and
\code{diesel-async} 0.9.2 (still 0.x; revisit only if connection acquisition ever becomes a
measured bottleneck).

\begin{keyidea}
Where this document and the code disagree, the code is right and this document is out of date. It
was written against commit \texttt{befcba5} plus the uncommitted work in progress, and the working
tree moved three times during the writing --- which is a good sign about the project and a caution
about the document.
\end{keyidea}
